%------------------------------------------------------------------------------%
\documentclass[a4paper,12pt,fleqn]{article}

\usepackage{integracao_e_series}
\usepackage{systeme}

\ShowAnswers

%------------------------------------------------------------------------------%
\begin{document}

\makeHeader{Integrais e Séries}{Prova 2 -- Turma 6}{29/11/24}

% 01
%------------------------------------------------------------------------------%
\exercise{20}
Calcule o valor médio da função \(f(x)=x^2\ln(x)\) no intervalo \([1, 3]\)

\begin{answer}
  Calculando a primitiva de $f$ usando integral por partes
  \(\int u\,dv = uv - \int v\,du\)
  \begin{align*}
    F(x) & = \int f(x) dx \\[2mm]
         & = \int x^2\ln(x)\, dx
  \end{align*}
  \[
    u = \ln(x)
    \qquad
    du = \frac{dx}{x}
    \qquad\qquad
    dv = x^2\, dx
    \qquad
    v = \frac{x^3}{3}
  \]
  \begin{align*}
    F(x) & = \frac{x^3}{3}\ln(x) - \int \frac{x^3}{3} \frac{dx}{x}\\[2mm]
         & = \frac{x^3}{3}\ln(x) - \frac{1}{3} \int x^2 \,dx\\[2mm]
         & = \frac{x^3}{3}\ln(x) - \frac{1}{3} \frac{x^3}{3} + C \\[2mm]
         & = \frac{x^3}{3}\left(\ln(x) -\frac{1}{3}\right) + C
  \end{align*}

  Calculando a integral definida de $f$ no intervalo \([1, 3]\)
  \begin{align*}
    I & = \int_1^3 f(x) dx   \\[2mm]
      & = F(x) \evalat{1}{3} \\[2mm]
      & = \left[\frac{x^3}{3}\left(\ln(x) -\frac{1}{3}\right)\right] \evalat{1}{3} \\[2mm]
      & = \frac{3^3}{3}\left(\ln(3) -\frac{1}{3}\right)
        - \frac{1^3}{3}\left(\ln(1) -\frac{1}{3}\right) \\[2mm]
      & = 9\left(\ln(3) -\frac{1}{3}\right)
        - \frac{1}{3}\left(0 -\frac{1}{3}\right) \\[2mm]
      & = 9\ln(3) -3 + \frac{1}{9}   \\[2mm]
      & = 9\ln(3) -\frac{27}{9} + \frac{1}{9}  \\[2mm]
      & = 9\ln(3) -\frac{26}{9}
  \end{align*}

  Calculando a média
  \begin{align*}
    \fmedio & = \frac{1}{b-a} \int_a^b f(x)dx \\[2mm]
            & = \frac{1}{3-1} \int_1^3 x^2\ln(x)\,dx \\[2mm]
            & = \frac{1}{2} \left(9\ln(3) -\frac{26}{9}\right) \\[2mm]
            & = \frac{9}{2}\ln(3) -\frac{13}{9} \\[2mm]
            & \approx 3,4993
  \end{align*}
\end{answer}

% 02
%------------------------------------------------------------------------------%
\exercise{20}
Use substituição simples para calcular a integral
\(
  \int_0^\infty x e^{-x^2}\,dx
\)

\begin{answer}
  Calculando a primitiva usando substituição simples
  \[
    F = \int x e^{-x^2}\,dx
  \]
  \[
    u = -x^2
    \qquad
    du = -2x\,dx
    \qquad
    x\,dx = \frac{du}{-2}
  \]
  \begin{align*}
    F & = \int e^{u}\,\frac{du}{-2} \\[2mm]
      & = -\frac{1}{2}\int e^{u}\,du \\[2mm]
      & = -\frac{1}{2} e^{u} + C \\[2mm]
      & = -\frac{1}{2} e^{-x^2} + C
  \end{align*}

  Calculando a integral imprópria
  \begin{align*}
    I & = \int_0^\infty x e^{-x^2}\,dx \\[2mm]
      & = \lim_{b\to\infty}\int_0^b x e^{-x^2}\,dx \\[2mm]
      & = \lim_{b\to\infty}\left(F(x)\evalat{0}{b}\right) \\[2mm]
      & = \lim_{b\to\infty}\left(-\frac{e^{-x^2}}{2} \evalat{0}{b}\right) \\[2mm]
      & = \lim_{b\to\infty}\left(-\frac{e^{-b^2}}{2} + \frac{e^{-0^2}}{2}  \right) \\[2mm]
      & = \lim_{b\to\infty}\left(-\frac{e^{-b^2}}{2} + \frac{1}{2} \right) \\[2mm]
      & = -\frac{1}{2}\lim_{b\to\infty}\left(e^{-b^2} \right) + \frac{1}{2} \\[2mm]
      & = \frac{1}{2}
  \end{align*}
\end{answer}

% 03
%------------------------------------------------------------------------------%
\exercise{20}
Encontre a primitiva da função \(f(x) = \cos^5(x)\sin^2(x) \)

\begin{answer}
  \begin{align*}
    F
    & = \int f(x) dx \\
    & = \int \cos^5(x)\sin^2(x) dx \\
    & = \int \left(\cos^2(x)\right)^2\sin^2(x)\cos(x) dx \\
    & = \int \left(1-\sin^2(x)\right)^2\sin^2(x)\cos(x) dx
  \end{align*}
  Fazendo a substituição
  \(
  \qquad
  u = \sin(x)
  \qquad
  du = \cos(x)dx
  \)
  \begin{align*}
    F
    & = \int \left(1-u^2\right)^2 u^2 du \\
    & = \int \left(1-2u^2+u^4\right) u^2 du \\
    & = \int u^2-2u^4+u^6 \, du \\
    & = \frac{1}{3} u^3
      - \frac{2}{5} u^5
      + \frac{1}{7} u^7
      + c \\
    & = \frac{1}{3} \sin^3(x)
      - \frac{2}{5} \sin^5(x)
      + \frac{1}{7} \sin^7(x)
      + c
  \end{align*}
\end{answer}


% 04
%------------------------------------------------------------------------------%
\exercise{20}
Encontre a primitiva da função \( h(x) = \frac{x^5}{\sqrt{4-x^2}} \)

\begin{answer}
  \[
  H = \int \frac{x^5}{\sqrt{4-x^2}} dx
  \]
  Fazendo a substituição
  \(
  x = 2\sin(\theta)
  \qquad
  dx = 2\cos(\theta) d\theta
  \)
  \begin{spreadlines}{2ex}
    \begin{align*}
      H
      & = \int \frac{2^5 \sin^5(\theta)}{\sqrt{4-4\sin^2(\theta)}} 2\cos(\theta) d\theta \\
      & = 2^5 \int \frac{\sin^5(\theta)}{\sqrt{1-\sin^2(\theta)}} \cos(\theta) d\theta \\
      & = 2^5 \int \frac{\sin^5(\theta)}{\cos(\theta)} \cos(\theta) d\theta \\
      & = 2^5 \int \sin^5(\theta) d\theta \\
      & = 2^5 \int \left(\sin^2(\theta)\right)^2 \sin(\theta) d\theta \\
      & = 2^5 \int \left(1-\cos^2(\theta)\right)^2 \sin(\theta) d\theta
    \end{align*}
  \end{spreadlines}
  Fazendo a substituição
  \(
  \qquad
  u = \cos(\theta)
  \qquad
  du = -\sin(\theta)d\theta
  \)
  \begin{align*}
    H
    & = -2^5 \int \left(1-u^2\right)^2 du \\
    & = -2^5 \int 1-2u^2 + u^4\, du \\
    & = 2^5 \int 2u^2 - u^4 -1\, du \\
    & = 2^5 \left(\frac{2}{3}u^3 - \frac{1}{5}u^5 - u \right) + c \\
    & = 2^5 \left(\frac{2}{3}\cos^3(\theta)
    - \frac{1}{5}\cos^5(\theta)
    - \cos(\theta) \right) + c
  \end{align*}
  Para calcular \(\cos(\theta)\) usamos que \(\sin(\theta) = \sfrac{x}{2}\), portanto
  a hipotenusa é 2 e o cateto oposto é $x$ assim o cateto adjacente é
  \(
  a = \sqrt{2^2-x^2} = \sqrt{4-x^2}
  \)
  e o cosseno é
  \[
  \cos(\theta) = \frac{\sqrt{4-x^2}}{2}
  \]
  Voltando para a integral
  \begin{spreadlines}{2ex}
    \begin{align*}
      H
      & = 2^5\left[
      \frac{2}{3}\left(\frac{\sqrt{4-x^2}}{2}\right)^3
      - \frac{1}{5}\left(\frac{\sqrt{4-x^2}}{2}\right)^5
      - \frac{\sqrt{4-x^2}}{2}
      \right] + c \\
      & = 2^5\left[
      \frac{2}{3\cdot 2^3}\left(\sqrt{4-x^2}\right)^3
      - \frac{1}{5\cdot 2^5}\left(\sqrt{4-x^2}\right)^5
      - \frac{1}{2} \sqrt{4-x^2}
      \right] + c \\
      & = \sqrt{4-x^2}
      \left[
      \frac{2^3}{3}\left(\sqrt{4-x^2}\right)^2
      - \frac{1}  {5}\left(\sqrt{4-x^2}\right)^4
      - 2^4
      \right] + c \\
      & = \sqrt{4-x^2}
      \left[
      \frac{8}{3}\left(4-x^2\right)
      - \frac{1}{5}\left(4-x^2\right)^2
      - 16
      \right] + c \\
      & = \sqrt{4-x^2}
      \left[
      \frac{32-8x^2}{3}
      - \frac{16-8x^2+x^4}{5}
      - 16
      \right] + c \\
      & = \sqrt{4-x^2}
      \left[
      \frac{5\cdot 32 - 5\cdot 8x^2}{15}
      -\frac{3\cdot 16 - 3\cdot 8x^2 + 3x^4}{15} - \frac{16\cdot 15}{15}
      \right] + c \\
      & = \frac{\sqrt{4-x^2}}{15}
      \Big[
      16\cdot 5\cdot 2 - 16 \cdot 3 - 16\cdot 5 \cdot 3
      - 5\cdot 8x^2 + 3\cdot 8x^2
      - 3x^4
      \Big] + c \\
      & = \frac{\sqrt{4-x^2}}{15}
      \Big[
      - 16\cdot 5 -16 \cdot 3
      - 2\cdot 8x^2
      - 3x^4
      \Big] + c \\
      & = -\frac{\sqrt{4-x^2}}{15}
      \Big[
      16\cdot 8
      + 2\cdot 8x^2
      + 3x^4
      \Big] + c \\
      & = -\frac{\sqrt{4-x^2}}{15}
      \Big[ 128 + 16x^2 + 3x^4 \Big] + c
    \end{align*}
  \end{spreadlines}
\end{answer}

% 05
%------------------------------------------------------------------------------%
\exercise{20}
Calcule o volume do sólido de revolução gerado pela rotação,
em torno da reta \(x=-2\), da região contida entre as curvas
\[
  f(x) = \frac{1}{5x-x^2} \qquad
  y = 0 \qquad
  x = 1 \qquad
  x = 4
\]

\begin{answer}
  Volume por castas cilíndricas
  \begin{align*}
    V & = \int_a^b 2\pi r(x) h(x) dx \\
    a & = 1 \\
    b & = 4 \\
    r & = x+2 \\
    h & = \frac{1}{5x-x^2}
  \end{align*}
  portanto
  \[
    V = 2\pi \int_1^4 \frac{x+2}{5x-x^2} dx
  \]
  Calculando a integral indefinida
  \[
    F = \int \frac{x+2}{5x-x^2} dx = \int \frac{x+2}{x(5-x)} dx
  \]
  Por frações parciais temos
  \begin{align*}
    \frac{x+2}{x(5-x)} & = \frac{A}{x} + \frac{B}{5-x} \\
    x+2                & = A(5-x) + Bx                 \\
    & = (B-A)x + 5A
  \end{align*}
  Igualando os coeficientes
  \[
    \systeme[AB]{B-A=1,5A=2}
  \]
  obtemos os valores
  \(A=\frac{2}{5}\) e
  \(B=\frac{7}{5}\)
  portanto
  \begin{align*}
    F & = \int \frac{2}{5x} + \frac{7}{5(5-x)} dx \\
      & = \frac{2}{5}\ln(x) - \frac{7}{5}\ln(5-x) + c
  \end{align*}
  Voltando ao volume temos
  \begin{align*}
    V
    & = 2\pi F(x)\evalat{1}{4} \\
    & = 2\pi \left[\frac{2}{5}\ln(x) - \frac{7}{5}\ln(5-x)\right] \evalat{1}{4} \\
    & = 2\pi \left[
      \left(\frac{2}{5}\ln(4) - \frac{7}{5}\ln(5-4)\right)
    - \left(\frac{2}{5}\ln(1) - \frac{7}{5}\ln(5-1)\right)
    \right] \\
    & = 2\pi \left(\frac{2}{5}\ln(4) + \frac{7}{5}\ln(4)\right) \\
    & = 2\pi \frac{9}{5}\ln(4) \\
    & = \frac{18\pi}{5} \ln(4)
  \end{align*}
\end{answer}

%------------------------------------------------------------------------------%
\end{document}
%------------------------------------------------------------------------------%
